\chapter{EDA II: Bivariate and Multivariate Analysis, Correlations and Multicollinearity}
\companion{StatQuest: Pearson's Correlation, Clearly Explained}{\ytid{xZ_z8KWkhXE}}

Univariate analysis describes one column. Bivariate analysis asks how \emph{two} variables move together (``does MBA percentage relate to placement?''),
and multivariate analysis asks how many variables relate at once (``are experience, age and graduation year telling the same story?''). The checklist
names Pearson, Spearman (with ranks), point-biserial and phi correlations, reading a correlation matrix, multicollinearity and the VIF formula in terms of
$R$. Each is derived here with a worked example.

\section{Which measure for which pair of variable types?}
\begin{center}\small
\begin{tabularx}{\linewidth}{l l X}
\toprule
Variable 1 & Variable 2 & Measure\\
\midrule
numeric & numeric & Pearson $r$ (linear), Spearman $\rho$ or Kendall $\tau$ (monotonic, robust)\\
numeric & binary & point-biserial correlation (a special case of Pearson)\\
numeric & multi-class categorical & ANOVA $F$-test, correlation ratio $\eta^2$, box plots\\
binary & binary & phi coefficient $\phi$; odds ratio\\
categorical & categorical & chi-square test, Cram\'er's $V$\\
ordinal & ordinal / numeric & Spearman, Kendall\\
\bottomrule
\end{tabularx}
\end{center}

\section{Covariance}
Covariance measures whether two variables move together:
\[ \Cov(x,y) = \frac{1}{n-1}\sum_i(x_i - \bar x)(y_i - \bar y). \]
If above-average $x$ tends to come with above-average $y$, the products are mostly positive: positive covariance. Its size depends on units (rupees vs lakh),
so it is hard to interpret alone; correlation fixes this.

\section{Pearson correlation}
\[ r = \frac{\Cov(x,y)}{s_xs_y} = \frac{\sum(x_i - \bar x)(y_i - \bar y)}{\sqrt{\sum(x_i - \bar x)^2}\sqrt{\sum(y_i - \bar y)^2}}. \]
It is the covariance of the standardised variables: unit-free, between $-1$ and $+1$.
\begin{example}[: Pearson by hand]
Years of experience $x = 1,2,3,4,5$; rating $y = 2,4,5,4,5$. $\bar x = 3$, $\bar y = 4$.
Deviations $x$: $-2,-1,0,1,2$; $y$: $-2,0,1,0,1$. Products: $4, 0, 0, 0, 2$; sum $= 6$.
$\sum(x - \bar x)^2 = 10$, $\sum(y - \bar y)^2 = 4 + 0 + 1 + 0 + 1 = 6$.
$r = 6/\sqrt{10\cdot6} = 6/7.746 = 0.775$. Covariance $= 6/4 = 1.5$.
\end{example}
\textbf{Properties and cautions}
\begin{itemize}
\item $r = \pm1$: points exactly on a line; $r = 0$: no \emph{linear} relationship (there may be a strong curved one: $y = x^2$ on symmetric $x$ has $r = 0$).
\item Unchanged by shifting or positive scaling of either variable (rupees or lakh: same $r$).
\item Sensitive to outliers: one extreme point can create or destroy a correlation.
\item $r^2$ is the fraction of variance of $y$ explained by a linear fit on $x$ (Chapter 2).
\item Correlation is not causation: ice-cream sales and drownings correlate through temperature (a confounder).
\end{itemize}

\section{Spearman rank correlation}
Replace each variable's values by their \textbf{ranks}, then compute Pearson on the ranks. It measures \textbf{monotonic} association (does $y$ tend to rise when
$x$ rises, in any shape?) and is robust to outliers and to non-linear but monotone relationships. Without ties, a shortcut:
\[ \rho = 1 - \frac{6\sum d_i^2}{n(n^2 - 1)},\qquad d_i = \text{rank}(x_i) - \text{rank}(y_i). \]
\begin{example}[: Spearman by hand (no ties)]
Interview scores $a = 10, 20, 30, 40, 50$ (ranks 1--5); manager ratings $b = 1, 3, 2, 5, 4$ (ranks 1, 3, 2, 5, 4).
$d = 0, -1, 1, -1, 1$; $\sum d^2 = 4$. $\rho = 1 - \frac{6\cdot4}{5\cdot24} = 1 - 0.2 = 0.8$.
\end{example}
\begin{example}[: ties]
For $y = 2, 4, 5, 4, 5$ the two 4s share ranks 2 and 3: both get 2.5; the two 5s get 4.5. Ranks: $1, 2.5, 4.5, 2.5, 4.5$. With ties, compute Pearson on these ranks
(the shortcut is only approximate): $\rho = 0.738$ against $x = 1,\dots,5$, slightly lower than Pearson's 0.775 here.
\end{example}
\textbf{Kendall's $\tau$}: based on concordant vs discordant pairs; more robust with small samples and many ties.

\section{Point-biserial correlation (numeric vs binary)}
Code the binary variable as 0/1 and compute Pearson; equivalently
\[ r_{pb} = \frac{\bar x_1 - \bar x_0}{s_x}\sqrt{p(1-p)}, \]
where $\bar x_1$, $\bar x_0$ are the means of the numeric variable in the two groups, $s_x$ its (population) standard deviation and $p$ the fraction in group 1.
\begin{example}[: point-biserial]
MBA \% $= 60, 65, 70, 75, 80, 85$; placed $= 0, 0, 1, 0, 1, 1$. $\bar x_1 = 78.33$, $\bar x_0 = 66.67$, $s_x = 8.54$, $p = 0.5$.
$r_{pb} = \frac{11.67}{8.54}\times0.5 = 0.683$. Placed students have clearly higher MBA scores.
\end{example}
It is the natural measure of how a numeric feature relates to a binary target; a $t$-test between the two groups tests the same thing.

\section{Phi coefficient (binary vs binary)}
For a $2\times2$ table $\begin{pmatrix}a & b\\c & d\end{pmatrix}$:
\[ \phi = \frac{ad - bc}{\sqrt{(a+b)(c+d)(a+c)(b+d)}}. \]
\begin{example}[: phi]
Rows: work experience yes/no; columns: placed yes/no. $a = 30$ (exp, placed), $b = 10$ (exp, not), $c = 20$ (no exp, placed), $d = 40$ (no exp, not).
$\phi = \frac{1200 - 200}{\sqrt{40\cdot60\cdot50\cdot50}} = \frac{1000}{2449.5} = 0.408$. Also $\phi = \sqrt{\chi^2/n} = \sqrt{16.67/100} = 0.408$.
\end{example}
For larger tables, \textbf{Cram\'er's $V$} $= \sqrt{\chi^2/(n(k - 1))}$ with $k = \min(\text{rows}, \text{cols})$, between 0 and 1.

\section{Numeric vs multi-class categorical: ANOVA and $\eta^2$}
Compare the means of a numeric variable across several groups (MBA \% by degree type). The \textbf{correlation ratio} $\eta^2 = \frac{\text{between-group sum of squares}}
{\text{total sum of squares}}$ is the fraction of variance explained by the group; the ANOVA $F$-test asks whether the differences are larger than chance (Chapter 39).

\section{Reading a correlation matrix}
A correlation matrix lists the correlation of every pair of numeric variables; usually shown as a heatmap.
\begin{center}\small
\begin{tabular}{l|ccccc}
& ssc & hsc & degree & mba & placed\\\hline
ssc & 1.00 & 0.51 & 0.54 & 0.39 & 0.61\\
hsc & 0.51 & 1.00 & 0.43 & 0.35 & 0.49\\
degree & 0.54 & 0.43 & 1.00 & 0.40 & 0.48\\
mba & 0.39 & 0.35 & 0.40 & 1.00 & 0.08\\
placed & 0.61 & 0.49 & 0.48 & 0.08 & 1.00\\
\end{tabular}
\end{center}
(Illustrative values.) How to read it, in order:
\begin{enumerate}
\item The diagonal is 1; the matrix is symmetric, so read one triangle.
\item \textbf{Target row/column} first: ssc has the strongest linear association with placement (0.61); mba almost none (0.08), surprisingly.
\item \textbf{Feature--feature} cells: high values ($|r| > 0.8$) warn of multicollinearity or redundancy; here moderate (0.4--0.54): the academic scores share a
  common ``academic ability'' factor.
\item \textbf{Blocks} of mutually correlated features suggest groups (PCA would find one component for them).
\item \textbf{Caveats}: only linear relationships (use Spearman for monotone), sensitive to outliers, hides non-linear and interaction effects, and with many variables
  some correlations are large by chance. Low correlation with the target does not make a feature useless (it may matter in interaction or non-linearly).
\end{enumerate}

\section{Simpson's paradox: when aggregate correlation lies}
A trend can reverse when data are split by a third variable. Overall, more certifications might correlate with lower salary; within each experience band the
correlation is positive, because juniors collect more certificates. Always check important relationships within key segments.

\section{Multicollinearity}
\subsection{What and why it matters}
Features are \textbf{multicollinear} when one can be (nearly) predicted from the others: age, years of experience and graduation year. Effects:
\begin{itemize}
\item Linear/logistic regression coefficients become unstable, with huge standard errors and sometimes the wrong sign; interpretation breaks.
\item Predictions are usually still fine, as long as the collinearity pattern holds in new data.
\item Tree models are unaffected in accuracy, but importance is split among the correlated features.
\end{itemize}
Pairwise correlations can miss it: $x_3 = x_1 + x_2$ may have only moderate pairwise correlations yet be perfectly collinear.

\subsection{VIF, derived}
Regress feature $j$ on all the other features and record its $R_j^2$ (how well the others explain it). Then
\[ \text{VIF}_j = \frac{1}{1 - R_j^2},\qquad \text{tolerance}_j = 1 - R_j^2. \]
The variance of the coefficient $\hat w_j$ is multiplied by VIF$_j$ compared with uncorrelated features; its standard error by $\sqrt{\text{VIF}_j}$.
\textbf{With only two predictors}, $R_j^2$ is just the squared correlation between them: VIF $= \frac{1}{1 - r^2}$.
\begin{example}[: VIF values]
Two predictors with $r = 0.9$: VIF $= 1/(1 - 0.81) = 5.26$; standard errors $2.3\times$ larger. $R_j^2 = 0.9$ gives VIF 10; $R_j^2 = 0.99$ gives 100.
\end{example}
Rules of thumb: VIF $>5$ worth a look, $>10$ serious. \textbf{Condition number} of $X$ (ratio of largest to smallest singular value, Chapter 14) $>30$ also signals it.

\subsection{Fixes}
Drop one of a correlated group (keep the most interpretable); combine them (an index, or ``experience $-$ age''); PCA; Ridge regression; collect data where they vary
independently; for prediction-only tasks, often leave it.

\subsection{The dummy-variable trap is perfect multicollinearity}
One-hot encoding a $k$-level category gives $k$ columns that always sum to 1, the same as the intercept column: perfect collinearity, $X^\top X$ singular
(Chapter 34). Drop one dummy (\code{drop\_first=True}) when the model has an intercept and no regularisation.

\begin{practical}[: where these measures are used in ML]
\begin{itemize}
\item \textbf{Feature screening}: rank features by $|r|$ or point-biserial correlation with the target as a first pass (then confirm with model-based importance).
\item \textbf{Redundancy removal}: drop one of each pair with $|r| > 0.95$ before linear models; check VIF.
\item \textbf{Categorical features}: Cram\'er's $V$ between candidate features (city and state are nearly the same information).
\item \textbf{Monitoring}: a change in the correlation between a feature and the target after deployment signals concept drift.
\item \textbf{Leakage detection}: a single feature with $|r| \approx 1$ with the target is suspicious (the salary-placement case).
\end{itemize}
\end{practical}

\stuck{StatQuest: Pearson's Correlation, Clearly Explained}{\ytid{xZ_z8KWkhXE}}
\section{Interview questions}

\begin{qa}{Compute the Pearson correlation of $x = 1,2,3,4,5$ and $y = 2,4,5,4,5$, and interpret it.}
\oneline{$r = 6/\sqrt{10\times6} = 0.775$: a fairly strong positive linear association; about $0.775^2 = 60\%$ of the variance in $y$ is explained by a straight line in $x$.}
\end{qa}

\begin{qa}{Pearson vs Spearman: what does each measure, and when do you prefer Spearman?}
\oneline{Pearson measures linear association of the raw values; Spearman is Pearson on ranks and measures monotonic association; prefer Spearman for ordinal data, non-linear
but monotone relationships, and data with outliers or heavy skew.}
\worked $y = e^x$ on $x = 1..5$: Spearman exactly 1 (perfectly monotone), Pearson about 0.89 (not linear).
\end{qa}

\begin{qa}{Write the Spearman formula and compute it for ranks $(1,2,3,4,5)$ and $(1,3,2,5,4)$.}
\oneline{$\rho = 1 - 6\sum d^2/(n(n^2 - 1)) = 1 - 6\cdot4/120 = 0.8$.}
\pushes \emph{``With ties?''} Assign average ranks and compute Pearson on ranks.
\end{qa}

\begin{qa}{How do you measure the association between a numeric feature and a binary target? Compute it for the MBA example.}
\oneline{With the point-biserial correlation, $r_{pb} = \frac{\bar x_1 - \bar x_0}{s_x}\sqrt{p(1-p)}$; for the example it is $0.683$.}
\why It equals Pearson with the binary variable coded 0/1. Significance via a two-sample $t$-test; a model-free alternative is the AUC of the feature alone.
\end{qa}

\begin{qa}{What is the phi coefficient? Compute it for the $2\times2$ table $[[30,10],[20,40]]$.}
\oneline{The correlation between two binary variables, $\phi = (ad - bc)/\sqrt{(a+b)(c+d)(a+c)(b+d)} = 1000/2449.5 = 0.408$; it equals $\sqrt{\chi^2/n}$.}
\end{qa}

\begin{qa}{How would you measure association between two categorical variables with many levels, like city and job function?}
\oneline{Use the chi-square test of independence for significance and Cram\'er's $V = \sqrt{\chi^2/(n(k-1))}$ for strength.}
\end{qa}

\begin{qa}{\deep $r = 0$ between two variables. Are they independent?}
\oneline{Not necessarily: $r = 0$ only rules out a linear relationship; a strong non-linear dependence (like $y = x^2$ for $x$ symmetric around 0) can have zero correlation.
For jointly normal variables, zero correlation does imply independence.}
\why Use mutual information, Spearman (for monotone), or plots to detect non-linear dependence.
\end{qa}

\begin{qa}{How do you read a correlation matrix in an interview?}
\oneline{Ignore the diagonal, read one triangle, look first at the target column for candidate predictors, then at feature pairs with high $|r|$ for redundancy and
multicollinearity, look for blocks of related features, and remember it shows only linear, pairwise, possibly confounded relationships.}
\end{qa}

\begin{qa}{Define VIF, write it in terms of $R$, and compute it for two predictors with correlation 0.9.}
\oneline{VIF$_j = 1/(1 - R_j^2)$ where $R_j^2$ is from regressing feature $j$ on the others; with two predictors $R_j^2 = r^2$, so VIF $= 1/(1 - 0.81) = 5.26$.}
\why It is the factor by which collinearity inflates the variance of the coefficient estimate.
\end{qa}

\begin{qa}{\deep Why can pairwise correlations miss multicollinearity?}
\oneline{Because collinearity can involve three or more variables jointly; for example $x_3 = x_1 + x_2$ with independent $x_1$, $x_2$ gives pairwise correlations of only
about 0.71, yet $x_3$ is perfectly predictable from the other two (VIF infinite).}
\end{qa}

\begin{qa}{Does multicollinearity hurt prediction? Does it hurt random forests?}
\oneline{It mainly hurts coefficient interpretation and stability in linear models, not prediction accuracy (if the pattern persists); random forests predict fine but split
importance among correlated features, making each look less important.}
\end{qa}

\begin{qa}{What is Simpson's paradox? Give a job-portal example.}
\oneline{A relationship that holds in every subgroup reverses (or vanishes) when the groups are combined, because of a confounder; for example, recommendations might have a lower
apply rate overall but a higher rate within each experience band because they are shown more to juniors, who apply less.}
\end{qa}

\begin{qa}{\deep A feature has correlation 0.02 with the target. Should you drop it?}
\oneline{Not on that alone: correlation captures only linear marginal effects, and the feature may matter non-linearly or through interactions (for example only for one city);
check model-based importance, mutual information and segment-wise plots first.}
\end{qa}

\begin{qa}{Why is correlation not causation? How could you establish causation?}
\oneline{Because a correlation can come from a confounder, reverse causation, selection bias or chance; causation is established by randomised experiments (A/B tests) or, failing
that, by careful causal designs (controlling for confounders, instrumental variables, difference-in-differences, regression discontinuity).}
\end{qa}
